\documentclass[10pt,letterpaper]{article}
\usepackage[spanish,es-nodecimaldot]{babel}
\usepackage{fontspec}
\setmainfont{TeX Gyre Heros}
\setmonofont{Nimbus Mono PS}[Scale=0.86]
\usepackage[letterpaper,top=1.8cm,bottom=1.8cm,left=1.7cm,right=1.7cm]{geometry}
\setlength{\headheight}{21pt}
\usepackage[table]{xcolor}
\usepackage{graphicx,longtable,array,booktabs,fancyhdr,hyperref}
\usepackage{pdflscape,multicol,soul,tikz,titlesec,lastpage,needspace}
\usetikzlibrary{positioning}
\definecolor{accent}{HTML}{7C3AED}
\definecolor{accentblue}{HTML}{2563EB}
\definecolor{accentcyan}{HTML}{06B6D4}
\definecolor{ink}{HTML}{172033}
\definecolor{muted}{HTML}{667085}
\definecolor{softpurple}{HTML}{F3EFFF}
\definecolor{softblue}{HTML}{EEF5FF}
\definecolor{linecolor}{HTML}{D8E1F0}
\definecolor{entidadmark}{HTML}{FFE45C}
\definecolor{atributomark}{HTML}{FF8FA3}
\definecolor{relacionmark}{HTML}{8EE7A2}
\hypersetup{colorlinks=true,linkcolor=accent,urlcolor=accentblue,pdftitle={Diccionario de datos - Comercial La Estrella},pdfauthor={Alex Ricardo Castaneda Rodriguez}}
\AtBeginDocument{\color{ink}}
\pagestyle{fancy}\fancyhf{}
\fancyhead[L]{\color{accent}\bfseries Base de Datos 1 · Proyecto 1}
\fancyhead[R]{\color{accentblue}Diccionario de datos}
\fancyfoot[L]{\color{muted}\footnotesize Comercial La Estrella}
\fancyfoot[R]{\color{accent}\footnotesize\thepage/\pageref*{LastPage}}
\renewcommand{\headrulewidth}{.8pt}
\renewcommand{\headrule}{\hbox to\headwidth{\color{accentblue}\leaders\hrule height \headrulewidth\hfill}}
\setlength{\parindent}{0pt}\setlength{\parskip}{4pt}\renewcommand{\arraystretch}{1.3}
\setlength{\LTpre}{5pt}\setlength{\LTpost}{11pt}
\titleformat{\section}[hang]{\Large\bfseries\color{accentblue}}{\colorbox{accent}{\color{white}\strut\thesection}}{.7em}{}
\titlespacing*{\section}{0pt}{1.1em}{.4em}
\newcommand{\gradientrule}[1]{\begin{tikzpicture}\shade[left color=accent,right color=accentblue] (0,0) rectangle (\textwidth,#1);\end{tikzpicture}}
\newcommand{\tabla}[2]{\Needspace{16\baselineskip}\section{#1}\noindent\colorbox{softpurple}{\parbox{\dimexpr\textwidth-2\fboxsep\relax}{\color{muted}\small #2}}\par\smallskip}
\newenvironment{dic}{%
\small\arrayrulecolor{linecolor}\rowcolors{2}{softblue}{white}%
\begin{longtable}{>{\ttfamily\color{accent}}p{3.45cm}p{2.75cm}>{\centering\arraybackslash}p{1.05cm}>{\centering\arraybackslash\bfseries\color{accentblue}}p{1.45cm}p{7.25cm}}
\rowcolor{accentblue}\multicolumn{1}{p{3.45cm}}{\color{white}\bfseries Columna} & \color{white}\bfseries Tipo & \color{white}\bfseries Nulo & \color{white}\bfseries Clave & \color{white}\bfseries Restricción y descripción\\
\endfirsthead
\rowcolor{accentblue}\multicolumn{1}{p{3.45cm}}{\color{white}\bfseries Columna} & \color{white}\bfseries Tipo & \color{white}\bfseries Nulo & \color{white}\bfseries Clave & \color{white}\bfseries Restricción y descripción\\
\endhead}{\\\bottomrule\end{longtable}\rowcolors{2}{}{}}
\newcommand{\entidad}[1]{{\sethlcolor{entidadmark}\hl{#1}}}
\newcommand{\atributo}[1]{{\sethlcolor{atributomark}\hl{#1}}}
\newcommand{\relacion}[1]{{\sethlcolor{relacionmark}\hl{#1}}}
\begin{document}
\begin{titlepage}\thispagestyle{empty}
\begin{tikzpicture}[remember picture,overlay]
\shade[inner color=accent!15,outer color=white] ([xshift=-.4cm,yshift=-.8cm]current page.north east) circle (5.2cm);
\shade[inner color=accentblue!12,outer color=white] ([xshift=.4cm,yshift=.4cm]current page.south west) circle (4.6cm);
\shade[left color=accent,right color=accentblue] (current page.north west) rectangle ([yshift=-.34cm]current page.north east);
\shade[left color=accentblue,right color=accentcyan] ([yshift=.24cm]current page.south west) rectangle (current page.south east);
\end{tikzpicture}
\begin{center}\vspace*{.15cm}
\begin{tikzpicture}\node[circle,shade,inner color=white,outer color=accent!18,draw=accentblue,line width=1.2pt,inner sep=7pt]{\includegraphics[width=3.05cm]{assets/escudo.png}};\end{tikzpicture}\par\medskip
{\Large\bfseries\color{accent}UNIVERSIDAD DE SAN CARLOS DE GUATEMALA}\par\smallskip
{\large\color{accentblue}Facultad de Ingeniería · Escuela de Ciencias y Sistemas}\par\vfill
\begin{tikzpicture}\node[rectangle,rounded corners=11pt,shade,left color=accent,right color=accentblue,text=white,align=center,text width=.86\textwidth,inner sep=17pt]{{\Huge\bfseries DICCIONARIO DE DATOS}\par\medskip{\Large Diseño e implementación de una base de datos\\para el control de ventas}};\end{tikzpicture}\par\medskip
{\Large\bfseries\color{accentblue}Comercial La Estrella}\par\vfill
\begin{tikzpicture}[node distance=.45cm]
\node[rounded corners=7pt,fill=softpurple,draw=accent!30,minimum width=3.2cm,minimum height=1.45cm,align=center] (a) {{\LARGE\bfseries\color{accent}19}\\[-1pt]{\small TABLAS}};
\node[right=.45cm of a,rounded corners=7pt,fill=softblue,draw=accentblue!30,minimum width=3.2cm,minimum height=1.45cm,align=center] (b) {{\LARGE\bfseries\color{accentblue}68}\\[-1pt]{\small ATRIBUTOS}};
\node[right=.45cm of b,rounded corners=7pt,fill=accentcyan!8,draw=accentcyan!35,minimum width=3.2cm,minimum height=1.45cm,align=center] (c) {{\LARGE\bfseries\color{accentcyan}3FN}\\[-1pt]{\small NORMALIZACIÓN}};
\end{tikzpicture}\par\vfill
\begin{tabular}{rl}\textbf{Estudiante:}&Alex Ricardo Castañeda Rodríguez\\\textbf{Carné:}&202300476\\\textbf{Curso:}&Bases de Datos 1\\\textbf{Proyecto:}&1\\\textbf{Fecha:}&Septiembre de 2026\end{tabular}\par\vfill
{\small\color{muted}Guatemala, 2026}
\end{center}
\end{titlepage}

\begin{landscape}
\thispagestyle{empty}
\begin{tikzpicture}[remember picture,overlay]
\shade[left color=accent,right color=accentblue] (current page.north west) rectangle ([yshift=-.22cm]current page.north east);
\shade[inner color=accentblue!8,outer color=white] ([xshift=-.2cm,yshift=-.2cm]current page.south east) circle (3.8cm);
\end{tikzpicture}
{\small\bfseries\color{accent}BASES DE DATOS 1 · PROYECTO 1}\par\smallskip
{\Huge\bfseries Análisis del enunciado}\par
{\large\color{muted}Identificación visual de entidades, atributos y relaciones}\par\medskip

\noindent\colorbox{entidadmark}{\strut\bfseries ENTIDADES}\hspace{1em}
\colorbox{atributomark}{\strut\bfseries ATRIBUTOS}\hspace{1em}
\colorbox{relacionmark}{\strut\bfseries RELACIONES Y REGLAS}\par\medskip

\begingroup
\fontsize{8.2}{10.2}\selectfont
\setlength{\columnsep}{1cm}
\begin{multicols}{2}
La empresa comercial \entidad{Comercial La Estrella} opera una red de \entidad{tiendas} dedicadas a la venta de \entidad{productos} de consumo general. La empresa requiere diseñar e implementar una base de datos relacional que permita administrar la información básica de sus tiendas, empleados, clientes, productos y ventas, así como obtener reportes que apoyen el análisis de sus operaciones comerciales.

\medskip
La empresa opera en distintos \entidad{países}. \relacion{Cada país se identifica mediante un código y un nombre}. Dentro de cada país existen varios \entidad{departamentos} y cada departamento agrupa varios \entidad{municipios}. \relacion{Un departamento pertenece a un único país y un municipio pertenece a un único departamento}.

\medskip
Cada \entidad{tienda} se identifica mediante un \atributo{código único}. De cada tienda se debe registrar su \atributo{nombre, dirección, teléfono y municipio} donde se encuentra ubicada. Además, cada tienda pertenece a un \entidad{tipo de tienda}, por ejemplo: tienda de conveniencia, supermercado, tienda mayorista o tienda especializada. \relacion{Un tipo de tienda puede clasificar a varias tiendas, pero cada tienda pertenece a un único tipo}.

\medskip
La empresa cuenta con \entidad{empleados} asignados a sus tiendas. De cada empleado se debe registrar un \atributo{código, nombres, apellidos, correo electrónico, teléfono, fecha de contratación y tienda a la que pertenece}. Cada empleado ocupa un \entidad{cargo} dentro de la empresa, por ejemplo: gerente de tienda, supervisor, vendedor o cajero. \relacion{Un cargo puede estar asignado a varios empleados; sin embargo, un empleado tiene un único cargo y pertenece a una única tienda}.

\medskip
Los \entidad{clientes} pueden realizar compras en cualquiera de las tiendas de la empresa. De cada cliente se debe almacenar un \atributo{código, nombres, apellidos, tipo de identificación, número de identificación, teléfono, correo electrónico, dirección y municipio de residencia}. Cada cliente registra un único \entidad{tipo de identificación}, mientras que \relacion{un mismo tipo de identificación puede utilizarse para varios clientes}.

\columnbreak

\medskip
La empresa administra un \entidad{catálogo de productos}. Cada \entidad{producto} posee un \atributo{código, nombre, descripción, precio vigente de venta y existencia actual}. Los productos se clasifican por \entidad{categoría} y \entidad{marca}. Cada categoría y cada marca poseen un \atributo{código y un nombre}. \relacion{Una categoría puede contener varios productos y una marca puede estar asociada con varios productos; cada producto pertenece a una única categoría y a una única marca}.

\medskip
Las \entidad{ventas} se realizan en una tienda y son atendidas por un empleado. Cada venta debe registrar un \atributo{número único, fecha, cliente, tienda, empleado responsable y estado actual}. Los \entidad{estados de una venta} incluyen, como mínimo: registrada, pagada y anulada. \relacion{Un cliente puede realizar varias ventas; una tienda puede registrar muchas ventas; y un empleado puede atender varias ventas}.

\medskip
Cada venta puede incluir uno o varios \entidad{productos}. Por cada producto vendido se debe registrar la \atributo{cantidad, el precio unitario aplicado durante la venta y el subtotal}. \relacion{El precio unitario del detalle debe conservar el valor aplicado al momento de realizar la operación, aunque el precio vigente del producto cambie posteriormente}.

\medskip
Una venta puede registrarse con uno o varios \entidad{pagos}. Para cada pago se debe indicar el \atributo{método de pago utilizado y el monto pagado}. La empresa maneja distintos \entidad{métodos de pago}, tales como efectivo, tarjeta de débito, tarjeta de crédito o transferencia. \relacion{La suma de los pagos de una venta pagada debe corresponder al total de los detalles de la venta}.

\medskip
La empresa necesita generar reportes sobre las ventas realizadas, los productos más vendidos, las categorías y marcas con mayor facturación, el desempeño de los empleados, las ventas por tienda y tipo de tienda, las compras realizadas por los clientes y los métodos de pago utilizados.
\end{multicols}
\endgroup

\vfill
\color{muted}\footnotesize Anexo del Diccionario de datos · Comercial La Estrella
\end{landscape}
\clearpage

{\Huge\bfseries\color{accent}Índice del diccionario}\par\smallskip
{\large\color{muted}Estructura documentada de la base de datos Comercial La Estrella}\par\smallskip
\gradientrule{.08cm}\par\medskip
\noindent\colorbox{softpurple}{\parbox{\dimexpr\textwidth-2\fboxsep\relax}{\small\textbf{Convenciones:}\quad \textcolor{accent}{\textbf{PK}} clave primaria \quad \textcolor{accentblue}{\textbf{FK}} clave foránea \quad \textcolor{accentcyan}{\textbf{UQ}} unicidad \quad * restricción compuesta}}\par\medskip
\tableofcontents\clearpage

\tabla{PAIS}{Catálogo de países donde opera la empresa.}
\begin{dic}id\_pais&NUMBER(10)&No&PK&Identificador del país.\\nombre&VARCHAR2(80)&No&UQ&Nombre único del país.\end{dic}
\tabla{DEPARTAMENTO}{División territorial perteneciente a un país.}
\begin{dic}id\_departamento&NUMBER(10)&No&PK&Identificador del departamento.\\nombre&VARCHAR2(80)&No&UQ*&Nombre; único junto con id\_pais.\\id\_pais&NUMBER(10)&No&FK&Referencia a PAIS.\end{dic}
\tabla{MUNICIPIO}{Municipio perteneciente a un departamento.}
\begin{dic}id\_municipio&NUMBER(10)&No&PK&Identificador del municipio.\\nombre&VARCHAR2(80)&No&UQ*&Nombre; único junto con id\_departamento.\\id\_departamento&NUMBER(10)&No&FK&Referencia a DEPARTAMENTO.\end{dic}
\tabla{TIPO\_TIENDA}{Clasificación de establecimientos comerciales.}
\begin{dic}id\_tipo\_tienda&NUMBER(10)&No&PK&Identificador del tipo.\\nombre&VARCHAR2(80)&No&UQ&Nombre único.\end{dic}
\tabla{TIPO\_IDENTIFICACION}{Catálogo de documentos de identificación.}
\begin{dic}id\_tipo\_identificacion&NUMBER(10)&No&PK&Identificador del tipo.\\nombre&VARCHAR2(50)&No&UQ&Nombre único.\end{dic}
\tabla{CARGO}{Catálogo de cargos laborales.}
\begin{dic}id\_cargo&NUMBER(10)&No&PK&Identificador del cargo.\\nombre&VARCHAR2(80)&No&UQ&Nombre único.\end{dic}
\tabla{CATEGORIA}{Clasificación funcional de productos.}
\begin{dic}id\_categoria&NUMBER(10)&No&PK&Identificador de categoría.\\nombre&VARCHAR2(80)&No&UQ&Nombre único.\end{dic}
\tabla{MARCA}{Fabricante o marca comercial del producto.}
\begin{dic}id\_marca&NUMBER(10)&No&PK&Identificador de marca.\\nombre&VARCHAR2(80)&No&UQ&Nombre único.\end{dic}
\tabla{ESTADO\_VENTA}{Estado vigente de una venta.}
\begin{dic}id\_estado\_venta&NUMBER(10)&No&PK&Identificador del estado.\\nombre&VARCHAR2(30)&No&UQ&REGISTRADA, PAGADA o ANULADA.\end{dic}
\tabla{METODO\_PAGO}{Medios aceptados para recibir pagos.}
\begin{dic}id\_metodo\_pago&NUMBER(10)&No&PK&Identificador del método.\\nombre&VARCHAR2(80)&No&UQ&Nombre único.\end{dic}
\tabla{PERSONA}{Datos comunes sin duplicación para empleados y clientes.}
\begin{dic}id\_persona&NUMBER(10)&No&PK&Identificador de persona.\\nombres&VARCHAR2(100)&No&&Nombres.\\apellidos&VARCHAR2(100)&No&&Apellidos.\\telefono&VARCHAR2(25)&No&&Teléfono conservado como texto.\\correo&VARCHAR2(150)&Sí&UQ&Correo único cuando existe; Oracle permite varios NULL.\\direccion&VARCHAR2(200)&No&&Dirección postal.\\id\_municipio&NUMBER(10)&No&FK&Municipio de residencia.\end{dic}
\tabla{TIENDA}{Establecimiento donde se realizan ventas.}
\begin{dic}id\_tienda&NUMBER(10)&No&PK&Código de tienda.\\nombre&VARCHAR2(120)&No&UQ&Nombre único.\\direccion&VARCHAR2(200)&No&&Dirección física.\\telefono&VARCHAR2(25)&No&&Teléfono.\\id\_municipio&NUMBER(10)&No&FK&Ubicación municipal.\\id\_tipo\_tienda&NUMBER(10)&No&FK&Clasificación de tienda.\end{dic}
\tabla{EMPLEADO}{Rol laboral de una persona dentro de una tienda.}
\begin{dic}id\_empleado&NUMBER(10)&No&PK&Código de empleado.\\fecha\_contratacion&DATE&No&&Fecha de ingreso; se valida contra SYSDATE por consulta.\\id\_tienda&NUMBER(10)&No&FK/UQ*&Tienda; integra UQ (empleado, tienda).\\id\_cargo&NUMBER(10)&No&FK&Cargo desempeñado.\\id\_persona&NUMBER(10)&No&FK/UQ&Una persona puede ser un solo empleado.\end{dic}
\tabla{CLIENTE}{Rol comercial de una persona compradora.}
\begin{dic}id\_cliente&NUMBER(10)&No&PK&Código de cliente.\\id\_tipo\_identificacion&NUMBER(10)&No&FK/UQ*&Tipo; integra identificación única.\\numero\_identificacion&VARCHAR2(30)&No&UQ*&Número; único dentro de su tipo.\\id\_persona&NUMBER(10)&No&FK/UQ&Una persona puede ser un solo cliente.\end{dic}
\tabla{PRODUCTO}{Definición maestra independiente de precio y existencia por tienda.}
\begin{dic}id\_producto&NUMBER(10)&No&PK&Código de producto.\\nombre&VARCHAR2(120)&No&&Nombre comercial.\\descripcion&VARCHAR2(300)&No&&Descripción.\\id\_categoria&NUMBER(10)&No&FK&Categoría.\\id\_marca&NUMBER(10)&No&FK&Marca.\end{dic}
\tabla{CATALOGO\_PRODUCTO}{Oferta e inventario vigente de cada producto en cada tienda.}
\begin{dic}precio\_vigente&NUMBER(12,2)&No&&Mayor que cero.\\existencia\_actual&NUMBER(10)&No&&Mayor o igual a cero.\\id\_tienda&NUMBER(10)&No&PK/FK&Primera parte de PK; referencia TIENDA.\\id\_producto&NUMBER(10)&No&PK/FK&Segunda parte de PK; referencia PRODUCTO.\end{dic}
\tabla{VENTA}{Cabecera de la transacción comercial.}
\begin{dic}id\_venta&NUMBER(10)&No&PK&Número único de venta.\\fecha\_venta&DATE&No&&Fecha de operación.\\id\_tienda&NUMBER(10)&No&FK*&Con id\_empleado referencia la asignación del empleado.\\id\_empleado&NUMBER(10)&No&FK*&Empleado de la misma tienda.\\id\_cliente&NUMBER(10)&No&FK&Cliente comprador.\\id\_estado\_venta&NUMBER(10)&No&FK&Estado actual.\end{dic}
\tabla{DETALLE\_VENTA}{Renglón de producto vendido; conserva el precio histórico.}
\begin{dic}cantidad&NUMBER(10)&No&&Mayor que cero.\\precio\_unitario&NUMBER(12,2)&No&&Precio aplicado, mayor que cero.\\subtotal&NUMBER(12,2)&No&&Igual a cantidad por precio unitario.\\id\_venta&NUMBER(10)&No&PK/FK&Primera parte de PK.\\id\_producto&NUMBER(10)&No&PK/FK&Segunda parte de PK; impide repetir producto.\end{dic}
\tabla{PAGO}{Abono individual aplicado a una venta.}
\begin{dic}id\_pago&NUMBER(10)&No&PK&Identificador del pago.\\monto&NUMBER(12,2)&No&&Mayor que cero.\\id\_metodo\_pago&NUMBER(10)&No&FK&Método utilizado.\\id\_venta&NUMBER(10)&No&FK&Venta receptora.\end{dic}

\Needspace{8\baselineskip}
{\Large\bfseries\color{accent}Leyenda y alcance}\par\smallskip
\gradientrule{.06cm}\par\smallskip
\begin{tikzpicture}
\node[rounded corners=7pt,fill=softpurple,draw=accent!25,text width=.94\textwidth,inner sep=10pt]{%
\textcolor{accent}{\textbf{PK}}: llave primaria \quad
\textcolor{accentblue}{\textbf{FK}}: llave foránea \quad
\textcolor{accentcyan}{\textbf{UQ}}: unicidad \quad
\textbf{*}: restricción compuesta.\par\medskip
Todas las columnas marcadas ``No'' son \texttt{NOT NULL}. Las reglas que dependen de agregados o de la fecha actual se comprueban en \texttt{05\_validaciones.sql}, porque Oracle no permite expresarlas como \texttt{CHECK} declarativo entre filas o tablas.};
\end{tikzpicture}
\end{document}
